	\section{Clustering}

	When many materials are compared using multiple isotherm parameters
	($Q_{\max}$, $K_L$, $E_{\mathrm{aff}}$, $\sigma_E$, etc.),
	unsupervised clustering can group materials with similar adsorption
	behaviour.

	\paragraph{$k$-Nearest Neighbours (KNN).}
	KNN is a non-parametric classifier that assigns a new material to
	the class most common among its $k$ nearest neighbours in
	parameter space.

	\subparagraph{Feature space and standardisation.}
	Each adsorbent is represented as a point in $d$-dimensional
	parameter space,
	$\mathbf{x}_i = (x_{i,1},\, x_{i,2},\, \ldots,\, x_{i,d})$.
	A typical choice is
	$\mathbf{x} = (Q_{\max},\; \ln K_{\mathrm{aff}}^{\ast},\;
	\sigma_E)$, so that $d = 3$.
	Because the parameters have different units and magnitudes,
	each feature is standardised to zero mean and unit variance
	before computing distances:
	\begin{equation}
		\tilde{x}_{i,\ell}
		= \frac{x_{i,\ell} - \bar{x}_\ell}{s_\ell},
		\qquad
		\bar{x}_\ell = \frac{1}{N}\sum_{i=1}^{N} x_{i,\ell},
		\qquad
		s_\ell = \sqrt{\frac{1}{N-1}\sum_{i=1}^{N}
			(x_{i,\ell} - \bar{x}_\ell)^2}.
	\end{equation}

	\subparagraph{Distance metric.}
	The distance between two materials $i$ and $j$ in standardised
	parameter space is
	\begin{equation}
		d(\mathbf{x}_i, \mathbf{x}_j)
		= \sqrt{\sum_{\ell=1}^{d}
			\left(\frac{x_{i,\ell} - x_{j,\ell}}{s_\ell}\right)^{\!2}},
		\label{eq:KNN_distance}
	\end{equation}
	where $s_\ell$ is the standard deviation of parameter~$\ell$.
	Alternative metrics (Manhattan, Mahalanobis) can be used when
	correlations between parameters are important.

	\subparagraph{Classification rule.}
	Given a query point $\mathbf{x}_q$ (a new material to classify),
	the algorithm proceeds as follows:
	\begin{enumerate}
		\item Compute $d(\mathbf{x}_q, \mathbf{x}_i)$ for all $N$
		known materials.
		\item Sort distances and select the $k$ nearest neighbours.
		\item Assign $\mathbf{x}_q$ to the class that appears most
		frequently among the $k$ neighbours (majority vote):
		\begin{equation}
			\hat{y}(\mathbf{x}_q)
			= \arg\max_{c}\;
			\sum_{i \in \mathcal{N}_k(\mathbf{x}_q)}
			\mathbf{1}[y_i = c],
			\label{eq:KNN_vote}
		\end{equation}
		where $\mathcal{N}_k(\mathbf{x}_q)$ is the set of $k$
		nearest neighbours and $y_i$ is the class label of
		material~$i$.
	\end{enumerate}

	\subparagraph{Choice of $k$.}
	The hyperparameter $k$ controls the bias--variance trade-off:
	\begin{itemize}
		\item $k = 1$: the decision boundary follows the data
		closely (low bias, high variance); sensitive to noise.
		\item Large $k$: the boundary becomes smoother (higher bias,
		lower variance); may blur class boundaries.
	\end{itemize}
	In practice, $k$ is chosen by cross-validation. A common
	starting point is $k = \sqrt{N}$, rounded to the nearest odd
	integer to avoid ties.

	\subparagraph{Application to adsorbent classification.}
	In the context of adsorption isotherms, KNN can be used to:
	\begin{itemize}
		\item Classify a new adsorbent into performance categories
		(e.g.\ ``high capacity / low affinity'' vs.\ ``low capacity
		/ high affinity'') based on the universal descriptors
		$\{Q_{\max},\, K_{\mathrm{aff}}^{\ast},\, \sigma_E\}$.
		\item Identify the most similar reference materials in a
		database, assisting material selection for a target
		application.
		\item Detect outliers: materials whose $k$ nearest neighbours
		all belong to a different class may warrant re-examination.
	\end{itemize}

	Figure~\ref{fig:KNN_illustration} illustrates the KNN
	classification procedure in the two-dimensional projection
	$(Q_{\max},\, \ln K_{\mathrm{aff}}^{\ast})$.

	\begin{figure}[H]
		\centering
		\begin{tikzpicture}[
			every node/.style={font=\small},
			classA/.style={fill=blue!60, draw=blue!80, circle,
				minimum size=8pt, inner sep=0pt},
			classB/.style={fill=red!60, draw=red!80, regular polygon,
				regular polygon sides=4, minimum size=9pt,
				inner sep=0pt},
			classC/.style={fill=green!60, draw=green!80, regular polygon,
				regular polygon sides=3, minimum size=10pt,
				inner sep=0pt},
			query/.style={fill=yellow!80, draw=black, thick, star,
				star points=5, minimum size=12pt, inner sep=0pt},
			arr/.style={-{Stealth[length=4pt]}, thick},
			dashcirc/.style={draw=black!60, densely dashed, thick},
			]
			% --- Axes ---
			\draw[arr] (-0.3,0) -- (9.5,0)
			node[below] {$Q_{\max}$ (mg\,g$^{-1}$)};
			\draw[arr] (0,-0.3) -- (0,7.0)
			node[above, rotate=90, anchor=south]
			{$\ln K_{\mathrm{aff}}^{\ast}$};

			% --- Class A: high capacity, low affinity (blue circles) ---
			\node[classA] (a1) at (6.5, 1.5) {};
			\node[classA] (a2) at (7.2, 2.0) {};
			\node[classA] (a3) at (7.8, 1.2) {};
			\node[classA] (a4) at (6.8, 2.5) {};
			\node[classA] (a5) at (8.2, 1.8) {};
			\node[classA] (a6) at (7.5, 2.8) {};

			% --- Class B: low capacity, high affinity (red squares) ---
			\node[classB] (b1) at (1.5, 5.5) {};
			\node[classB] (b2) at (2.0, 6.0) {};
			\node[classB] (b3) at (2.5, 5.0) {};
			\node[classB] (b4) at (1.8, 4.5) {};
			\node[classB] (b5) at (2.8, 5.8) {};

			% --- Class C: moderate (green triangles) ---
			\node[classC] (c1) at (4.0, 3.5) {};
			\node[classC] (c2) at (4.5, 4.0) {};
			\node[classC] (c3) at (5.0, 3.0) {};
			\node[classC] (c4) at (4.2, 4.5) {};
			\node[classC] (c5) at (5.5, 3.8) {};
			\node[classC] (c6) at (3.5, 3.8) {};

			% --- Query point ---
			\node[query] (q) at (4.8, 3.5) {};
			\node[above right=2pt, font=\footnotesize\bfseries]
			at (q.north east) {$\mathbf{x}_q$};

			% --- KNN circle (k=5) ---
			\draw[dashcirc] (q) circle (1.6cm);
			\node[font=\footnotesize, anchor=west] at (6.0, 4.8)
			{$k = 5$ neighbourhood};
			\draw[thin, black!50, ->] (6.0, 4.7) -- (5.5, 4.1);

			% --- Distance lines to nearest neighbours ---
			\draw[thin, black!40, dotted] (q) -- (c2);
			\draw[thin, black!40, dotted] (q) -- (c3);
			\draw[thin, black!40, dotted] (q) -- (c5);
			\draw[thin, black!40, dotted] (q) -- (c1);
			\draw[thin, black!40, dotted] (q) -- (c4);

			% --- Legend ---
			\node[classA, label={[font=\footnotesize]right:
				High $Q_{\max}$, low $K_{\mathrm{aff}}^{\ast}$}]
			at (0.5, -1.0) {};
			\node[classB, label={[font=\footnotesize]right:
				Low $Q_{\max}$, high $K_{\mathrm{aff}}^{\ast}$}]
			at (4.0, -1.0) {};
			\node[classC, label={[font=\footnotesize]right:
				Moderate}]
			at (7.5, -1.0) {};
			\node[query, label={[font=\footnotesize]right:
				Query}]
			at (0.5, -1.7) {};

			% --- Annotation: vote result ---
			\node[draw, rounded corners, fill=green!10,
			font=\footnotesize, text width=3.2cm, align=center]
			at (8.0, 5.5)
			{Majority vote:\\
				5/5 = Moderate\\
				$\Rightarrow \hat{y} = $ \textbf{Moderate}};
		\end{tikzpicture}
		\caption{Illustration of $k$-Nearest Neighbours (KNN)
			classification in adsorbent parameter space.
			Each known material is plotted as a point in the
			$(Q_{\max},\, \ln K_{\mathrm{aff}}^{\ast})$ plane and
			coloured by its performance class.
			A query material~$\mathbf{x}_q$ (star) is classified by
			majority vote among its $k = 5$ nearest neighbours
			(dashed circle).  In this example, all five neighbours
			belong to the ``Moderate'' class (green triangles), so
			the query is assigned to that class.}
		\label{fig:KNN_illustration}
	\end{figure}

	\paragraph{Support Vector Machines (SVM).}
	SVM finds the optimal hyperplane that maximally separates classes
	in parameter space.

	\subparagraph{Linear SVM: maximum-margin hyperplane.}
	Given $N$ training materials with feature vectors
	$\mathbf{x}_i \in \mathbb{R}^d$ and class labels
	$y_i \in \{-1,+1\}$, SVM seeks the hyperplane
	$\mathbf{w}\cdot\mathbf{x} + b = 0$ that maximises the margin
	$2/\|\mathbf{w}\|$ between the two classes.
	The optimisation problem is
	\begin{equation}
		\min_{\mathbf{w},\,b}\;
		\frac{1}{2}\|\mathbf{w}\|^2
		\quad\text{subject to}\quad
		y_i\bigl(\mathbf{w}\cdot\mathbf{x}_i + b\bigr) \ge 1
		\quad \forall\, i.
		\label{eq:SVM_hard}
	\end{equation}
	The data points that lie exactly on the margin boundaries
	($y_i(\mathbf{w}\cdot\mathbf{x}_i + b) = 1$) are called
	\emph{support vectors}; they alone determine the decision
	boundary.

	\subparagraph{Soft-margin SVM.}
	When classes overlap (as is common in real adsorption data),
	slack variables $\xi_i \ge 0$ allow controlled misclassification:
	\begin{equation}
		\min_{\mathbf{w},\,b,\,\boldsymbol{\xi}}\;
		\frac{1}{2}\|\mathbf{w}\|^2
		+ C_{\mathrm{reg}}\sum_{i=1}^{N}\xi_i,
		\quad\text{s.t.}\quad
		y_i\bigl(\mathbf{w}\cdot\mathbf{x}_i + b\bigr)
		\ge 1 - \xi_i,
		\label{eq:SVM_soft}
	\end{equation}
	where $C_{\mathrm{reg}}$ is a regularisation parameter that
	controls the trade-off between a large margin and few
	misclassifications.

	\subparagraph{Kernel trick for non-linear boundaries.}
	For non-linearly separable data, a kernel function
	$\mathcal{K}(\mathbf{x}_i, \mathbf{x}_j)
	= \phi(\mathbf{x}_i)\cdot\phi(\mathbf{x}_j)$
	implicitly maps the data into a higher-dimensional space where
	a linear separator exists. Common kernels include:
	\begin{itemize}
		\item Linear: $\mathcal{K} = \mathbf{x}_i \cdot \mathbf{x}_j$.
		\item Radial basis function (RBF):
		$\mathcal{K} = \exp\!\bigl(-\gamma\|\mathbf{x}_i
		- \mathbf{x}_j\|^2\bigr)$.
		\item Polynomial:
		$\mathcal{K} = (\mathbf{x}_i\cdot\mathbf{x}_j + r)^p$.
	\end{itemize}
	The RBF kernel is the most widely used default, with
	$\gamma$ and $C_{\mathrm{reg}}$ selected by cross-validation.

	\subparagraph{Application to adsorbent classification.}
	SVM is particularly useful when:
	\begin{itemize}
		\item The number of materials is small relative to the number
		of descriptors ($N \lesssim 10\,d$), where SVM generalises
		better than KNN.
		\item A clear decision boundary between performance classes
		is desired (e.g.\ ``suitable'' vs.\ ``unsuitable'' for a
		given application).
		\item The feature space includes correlated parameters
		($Q_{\max}$, $K_{\mathrm{aff}}^{\ast}$, $\sigma_E$),
		where the kernel trick can capture non-linear interactions.
	\end{itemize}

	Figure~\ref{fig:SVM_illustration} illustrates the SVM
	classification procedure for a binary separation task in
	the $(Q_{\max},\, \ln K_{\mathrm{aff}}^{\ast})$ plane.

	\begin{figure}[H]
		\centering
		\begin{tikzpicture}[
			every node/.style={font=\small},
			classA/.style={fill=blue!60, draw=blue!80, circle,
				minimum size=8pt, inner sep=0pt},
			classB/.style={fill=red!60, draw=red!80, regular polygon,
				regular polygon sides=4, minimum size=9pt,
				inner sep=0pt},
			sv/.style={draw=black, very thick, minimum size=14pt,
				inner sep=0pt, fill opacity=0.6},
			arr/.style={-{Stealth[length=4pt]}, thick},
			]
			% --- Axes ---
			\draw[arr] (-0.3,0) -- (9.5,0)
			node[below] {$Q_{\max}$ (mg\,g$^{-1}$)};
			\draw[arr] (0,-0.3) -- (0,7.0)
			node[above, rotate=90, anchor=south]
			{$\ln K_{\mathrm{aff}}^{\ast}$};

			% --- Class A: high capacity, low affinity (blue circles) ---
			\node[classA] at (6.0, 1.2) {};
			\node[classA] at (7.0, 1.8) {};
			\node[classA] at (7.8, 1.0) {};
			\node[classA] at (6.5, 2.3) {};
			\node[classA] at (8.2, 1.5) {};
			\node[classA] at (7.5, 2.6) {};
			\node[classA] at (8.5, 2.0) {};
			% support vectors (class A)
			\node[classA, sv] (sva1) at (5.2, 2.8) {};
			\node[classA, sv] (sva2) at (6.0, 3.2) {};

			% --- Class B: low capacity, high affinity (red squares) ---
			\node[classB] at (1.5, 5.5) {};
			\node[classB] at (2.0, 6.2) {};
			\node[classB] at (2.5, 5.0) {};
			\node[classB] at (1.8, 4.8) {};
			\node[classB] at (3.0, 5.6) {};
			\node[classB] at (1.2, 5.8) {};
			% support vectors (class B)
			\node[classB, sv] (svb1) at (3.5, 4.2) {};
			\node[classB, sv] (svb2) at (2.8, 4.5) {};

			% --- Decision boundary (separating hyperplane) ---
			\draw[very thick, black] (1.0, 2.0) -- (8.5, 5.6)
			node[pos=0.15, above, rotate=27,
			font=\footnotesize\bfseries]
			{$\mathbf{w}\cdot\mathbf{x} + b = 0$};

			% --- Margin boundaries ---
			\draw[thick, blue!60, dashed] (1.0, 1.0) -- (8.5, 4.6);
			\draw[thick, red!60, dashed] (1.0, 3.0) -- (8.5, 6.6);

			% --- Margin annotation ---
			\draw[{Stealth[length=3pt]}-{Stealth[length=3pt]},
			thick, black!70]
			(4.3, 2.35) -- (4.3, 3.95)
			node[midway, left, font=\footnotesize]
			{margin $\frac{2}{\|\mathbf{w}\|}$};

			% --- Support vector labels ---
			\node[font=\tiny, anchor=west] at (6.2, 3.4)
			{support vectors};
			\draw[thin, black!50, ->] (6.2, 3.3) -- (6.05, 3.22);
			\draw[thin, black!50, ->] (6.2, 3.3) -- (5.35, 2.85);

			\node[font=\tiny, anchor=east] at (2.5, 4.0)
			{support vectors};
			\draw[thin, black!50, ->] (2.6, 4.05) -- (2.75, 4.42);
			\draw[thin, black!50, ->] (2.6, 4.05) -- (3.4, 4.18);

			% --- Legend ---
			\node[classA, label={[font=\footnotesize]right:
				Class~A (high $Q_{\max}$)}]
			at (0.5, -1.0) {};
			\node[classB, label={[font=\footnotesize]right:
				Class~B (high $K_{\mathrm{aff}}^{\ast}$)}]
			at (4.0, -1.0) {};
			\node[classA, sv, label={[font=\footnotesize]right:
				Support vectors}]
			at (8.0, -1.0) {};
		\end{tikzpicture}
		\caption{Illustration of Support Vector Machine (SVM)
			classification in adsorbent parameter space.
			The solid line is the maximum-margin hyperplane
			$\mathbf{w}\cdot\mathbf{x}+b=0$ that separates
			class~A (blue circles, high capacity) from class~B
			(red squares, high affinity).  Dashed lines mark the
			margin boundaries; the circled points on these
			boundaries are the support vectors that define the
			decision surface.  The margin width
			$2/\|\mathbf{w}\|$ is maximised by the SVM
			optimisation (Eq.~\ref{eq:SVM_hard}).}
		\label{fig:SVM_illustration}
	\end{figure}

	\subsection{Validation framework: incorporating uncertainty into
		machine-learning classification}

	The universal descriptors
	$\{Q_{\max},\, K_{\mathrm{aff}}^{\ast},\, E_{\mathrm{aff}},\,
	\sigma_E\}$ are not exact numbers: they are \emph{estimated}
	from finite, noisy experimental data via non-linear regression.
	Each descriptor therefore carries an associated standard
	deviation (or confidence interval) inherited from the fitting
	procedure.  Any classification analysis (KNN, SVM, or
	clustering) must account for this uncertainty; otherwise,
	materials whose descriptor values overlap within error may be
	incorrectly assigned to different classes.

	\subsubsection{Propagation of fitting uncertainty to descriptors}

	Let $\boldsymbol{\theta}$ be the vector of fitted isotherm
	parameters (e.g.\ $Q_{\max}$, $K_L$ for Langmuir) and
	$\boldsymbol{\Sigma}_{\theta}$ its covariance matrix, estimated
	from the Jacobian of the non-linear fit.  A universal descriptor
	$\phi = \phi(\boldsymbol{\theta})$ (for instance
	$K_{\mathrm{aff}} = K_H / Q_{\max}$) has variance
	\begin{equation}
		\sigma_\phi^2
		= \left(\frac{\partial\phi}{\partial\boldsymbol{\theta}}
		\right)^{\!\top}
		\boldsymbol{\Sigma}_{\theta}\,
		\frac{\partial\phi}{\partial\boldsymbol{\theta}},
		\label{eq:error_prop}
	\end{equation}
	obtained by first-order error propagation.  For each material
	the feature vector becomes
	$\mathbf{x}_i \pm \boldsymbol{\sigma}_i$, where
	$\boldsymbol{\sigma}_i
	= (\sigma_{Q_{\max}},\, \sigma_{\ln K^{\ast}},\,
	\sigma_{\sigma_E})_i$.

	\subsubsection{Incorporating standard deviations into KNN}

	The standard Euclidean metric
	(Eq.~\eqref{eq:KNN_distance}) treats all points as exact.
	When uncertainties are available, the distance can be
	modified to account for the ``size'' of each data point:

	\paragraph{Uncertainty-weighted distance.}
	\begin{equation}
		d_{\sigma}(\mathbf{x}_i, \mathbf{x}_j)
		= \sqrt{\sum_{\ell=1}^{d}
			\frac{(x_{i,\ell} - x_{j,\ell})^2}
			{\sigma_{i,\ell}^2 + \sigma_{j,\ell}^2}},
		\label{eq:KNN_sigma_distance}
	\end{equation}
	where $\sigma_{i,\ell}$ and $\sigma_{j,\ell}$ are the standard
	deviations of feature~$\ell$ for materials $i$ and $j$.
	This metric down-weights dimensions where both points have
	large uncertainty (the separation is not statistically
	significant) and up-weights dimensions where the data are
	precise.

	\paragraph{Probabilistic KNN via Monte Carlo.}
	An alternative approach samples $M$ realisations of each
	feature vector from its uncertainty distribution,
	$\mathbf{x}_i^{(m)} \sim \mathcal{N}(\mathbf{x}_i,\,
	\mathrm{diag}(\boldsymbol{\sigma}_i^2))$,
	and runs KNN on each realisation.  The class assignment
	becomes a probability:
	\begin{equation}
		P(\hat{y} = c \mid \mathbf{x}_q)
		= \frac{1}{M}\sum_{m=1}^{M}
		\mathbf{1}\!\left[
			\hat{y}^{(m)}(\mathbf{x}_q^{(m)}) = c
			\right].
		\label{eq:KNN_MC}
	\end{equation}
	A material is classified with high confidence only if
	$P(\hat{y} = c) \gg 0.5$; otherwise it lies in an
	\emph{uncertain region} between classes.

	\subsubsection{Incorporating standard deviations into SVM}

	\paragraph{Feature weighting.}
	Before training the SVM, features can be weighted inversely
	by their average uncertainty:
	\begin{equation}
		\tilde{x}_{i,\ell}
		= \frac{x_{i,\ell}}
		{\bar{\sigma}_\ell},
		\qquad
		\bar{\sigma}_\ell
		= \frac{1}{N}\sum_{i=1}^{N}\sigma_{i,\ell},
	\end{equation}
	so that noisy features contribute less to the margin
	calculation.  This is equivalent to using a weighted kernel
	$\mathcal{K}_\sigma(\mathbf{x}_i, \mathbf{x}_j)
	= \mathcal{K}(\tilde{\mathbf{x}}_i, \tilde{\mathbf{x}}_j)$.

	\paragraph{Fuzzy SVM.}
	In the soft-margin formulation
	(Eq.~\eqref{eq:SVM_soft}), each training point can be
	assigned a \emph{membership weight} $\mu_i \in (0,1]$
	inversely proportional to its total uncertainty:
	\begin{equation}
		\mu_i = \frac{1}{1 + \alpha\,\|\boldsymbol{\sigma}_i\|},
		\label{eq:fuzzy_weight}
	\end{equation}
	where $\alpha > 0$ is a scaling factor.  The modified
	objective becomes
	\begin{equation}
		\min_{\mathbf{w},\,b,\,\boldsymbol{\xi}}\;
		\frac{1}{2}\|\mathbf{w}\|^2
		+ C_{\mathrm{reg}}\sum_{i=1}^{N}\mu_i\,\xi_i,
		\label{eq:fuzzy_SVM}
	\end{equation}
	which reduces the penalty for misclassifying uncertain
	data points while maintaining a strict margin for
	well-determined materials.

	\paragraph{Monte Carlo SVM.}
	As with KNN, $M$ bootstrap samples can be drawn from the
	uncertainty distributions and an SVM trained on each.
	The ensemble of $M$ decision boundaries defines a
	\emph{margin uncertainty band} that quantifies how
	sensitive the classification is to measurement noise.

	\subsubsection{Cross-validation strategy}

	Both KNN and SVM require hyperparameter selection ($k$ for KNN;
	$C_{\mathrm{reg}}$ and $\gamma$ for SVM) and an estimate of
	generalisation performance.  The recommended protocol is:

	\begin{enumerate}
		\item \textbf{Stratified $k$-fold cross-validation.}
		Divide the $N$ materials into $k_{\mathrm{cv}}$ folds
		(typically $k_{\mathrm{cv}} = 5$ or $10$), preserving the
		class proportions in each fold.  For each fold,
		train on $k_{\mathrm{cv}}-1$ folds and evaluate on the
		held-out fold.

		\item \textbf{Nested cross-validation.}
		Use an outer loop for performance estimation and an inner
		loop for hyperparameter tuning, to avoid optimistic bias:
		\begin{equation}
			\text{Accuracy}
			= \frac{1}{k_{\mathrm{cv}}}
			\sum_{j=1}^{k_{\mathrm{cv}}}
			\frac{1}{|\mathcal{T}_j|}
			\sum_{i \in \mathcal{T}_j}
			\mathbf{1}[\hat{y}_i = y_i],
		\end{equation}
		where $\mathcal{T}_j$ is the test set in fold~$j$.

		\item \textbf{Uncertainty-aware scoring.}
		Beyond accuracy, report the \emph{classification
		uncertainty} for each material using the Monte Carlo
		approach (Eqs.~\eqref{eq:KNN_MC}) and flag materials
		with $P(\hat{y} = c) < 0.7$ as ``borderline''.

		\item \textbf{Reporting.}
		For each material, report:
		\begin{equation}
			\bigl\{
			Q_{\max} \pm \sigma_{Q},\;
			K_{\mathrm{aff}}^{\ast} \pm \sigma_{K},\;
			E_{\mathrm{aff}} \pm \sigma_{E_{\mathrm{aff}}},\;
			\sigma_E \pm \sigma_{\sigma_E},\;
			\hat{y},\;
			P(\hat{y})
			\bigr\},
		\end{equation}
		where the $\pm$ terms are standard deviations from the
		non-linear fit and $P(\hat{y})$ is the classification
		confidence from the Monte Carlo ensemble.
	\end{enumerate}

	\subsubsection{When to use KNN vs.\ SVM}

	Table~\ref{tab:KNN_vs_SVM} summarises the practical
	trade-offs between the two classifiers in the context of
	adsorbent comparison.

	\begin{table}[H]
		\centering
		\small
		\caption{Comparison of KNN and SVM for adsorbent
			classification based on universal descriptors.}
		\label{tab:KNN_vs_SVM}
		\renewcommand{\arraystretch}{1.4}
		\begin{tabular}{lcc}
			\hline
			\textbf{Criterion}
			& \textbf{KNN}
			& \textbf{SVM}
			\\
			\hline
			Training required
			& None (lazy learner)
			& Yes (optimisation)
			\\
			Decision boundary
			& Piecewise, local
			& Global hyperplane
			\\
			Sensitivity to noise
			& High (small $k$)
			& Low (margin-based)
			\\
			Handles uncertainty
			& Eq.~\eqref{eq:KNN_sigma_distance}
			& Eq.~\eqref{eq:fuzzy_SVM}
			\\
			Small datasets ($N < 30$)
			& Unstable
			& Preferred
			\\
			Large datasets ($N > 100$)
			& Robust
			& Excellent
			\\
			Interpretability
			& High (neighbours)
			& Moderate (support vectors)
			\\
			Multi-class
			& Native
			& One-vs-one / one-vs-rest
			\\
			\hline
		\end{tabular}
	\end{table}

	In practice, both methods should be applied in parallel.
	Agreement between KNN and SVM classifications increases
	confidence in the result; disagreement highlights materials
	that lie near class boundaries and warrant further
	experimental characterisation.

	% \subsection{Isotherm construction recommendations}
	% \subsubsection{Examples}
	% \subsection{Plot of variables}

	% \begin{itemize}
	% 	\item Number of points
	% \end{itemize}

